% Template Tese de Doutorado abnTeX2 — pesquisador-br-skill
% Suporta modalidade tradicional E "tese por artigos"
% Compilação: pdflatex + bibtex + pdflatex + pdflatex

\documentclass[
    12pt,
    openright,
    twoside,
    a4paper,
    chapter=TITLE,
    sumario=tradicional,
    section=TITLE,
    english,
    brazil
]{abntex2}

\usepackage{cmap}
\usepackage[utf8]{inputenc}
\usepackage[T1]{fontenc}
\usepackage{lmodern}
\usepackage{lastpage}
\usepackage{indentfirst}
\usepackage{color}
\usepackage{graphicx}
\usepackage{microtype}
\usepackage[brazilian,hyperpageref]{backref}
\usepackage[alf]{abntex2cite}
\usepackage{amsmath, amssymb}
\usepackage{booktabs}
\usepackage{pdfpages}               % pra incluir PDFs (artigos)

\titulo{TÍTULO DA TESE: subtítulo}
\autor{Nome Completo Sobrenome}
\local{Manaus -- AM}
\data{2026}
\orientador{Prof.\textsuperscript{a} Dr.\textsuperscript{a} Nome do Orientador}
% \coorientador{Prof. Dr. Nome do Coorientador}
\instituicao{Universidade Federal do Amazonas \par
            Programa de Pós-Graduação em [Nome] \par
            Doutorado em [Área]}
\tipotrabalho{Tese (Doutorado)}
\preambulo{Tese apresentada ao Programa de Pós-Graduação em
[Nome] da Universidade Federal do Amazonas, como requisito
parcial para a obtenção do título de Doutor em [Área].
\par
Linha de pesquisa: [Linha].}

\usepackage{hyperref}
\hypersetup{
    pdftitle={\@title},
    pdfauthor={\@author},
    colorlinks=true,
    linkcolor=blue,
    citecolor=blue,
    urlcolor=blue,
}

\begin{document}
\selectlanguage{brazil}
\frenchspacing

\pretextual
\imprimircapa
\imprimirfolhaderosto

% Folha de aprovação (5 examinadores: orientador + 2 internos + 2 externos)
\begin{folhadeaprovacao}
\begin{center}
{\ABNTEXchapterfont\large\imprimirautor}
\vspace*{\fill}
\begin{center}
\ABNTEXchapterfont\bfseries\Large\imprimirtitulo
\end{center}
\vspace*{\fill}

\hspace{.45\textwidth}\begin{minipage}{.5\textwidth}
\imprimirpreambulo
\end{minipage}\vspace*{\fill}
\end{center}

Tese aprovada. Manaus -- AM, \today:

\assinatura{\textbf{\imprimirorientadorRotulo} \\ \imprimirorientador}
\assinatura{\textbf{Prof. Dr. Nome 1 -- Examinador Interno} \\ Universidade Federal do Amazonas}
\assinatura{\textbf{Prof. Dr. Nome 2 -- Examinador Interno} \\ Universidade Federal do Amazonas}
\assinatura{\textbf{Prof. Dr. Nome 3 -- Examinador Externo} \\ [Outra Instituição]}
\assinatura{\textbf{Prof. Dr. Nome 4 -- Examinador Externo} \\ [Outra Instituição]}
\end{folhadeaprovacao}

\begin{dedicatoria}\vspace*{\fill}\centering\textit{Dedicatória.}\vspace*{\fill}\end{dedicatoria}

\begin{agradecimentos}
O presente trabalho foi realizado com apoio da [CAPES/CNPq/FAPESP]
-- processo nº [...].

Agradecimentos à orientação, à banca, ao programa, aos pares.
\end{agradecimentos}

\begin{epigrafe}\vspace*{\fill}\begin{flushright}
\textit{``Citação inspiradora.''} \\ (AUTOR, ANO, p. X)
\end{flushright}\end{epigrafe}

\setlength{\absparsep}{18pt}
\begin{resumo}
\noindent A presente tese investigou [problema] sob [N]
perspectivas articuladas, adotando [abordagem]. Os objetivos
foram: (i) [obj1]; (ii) [obj2]; (iii) [obj3]. A pesquisa foi
conduzida em [período], com [N] estudos articulados. Os achados
indicaram [resultado principal], contribuindo para [campo] ao
[contribuição]. As limitações incluem [limitação], indicando-se
[agenda].

\noindent\textbf{Palavras-chave}: termo 1; termo 2; termo 3;
termo 4; termo 5; termo 6.
\end{resumo}

\begin{resumo}[Abstract]\begin{otherlanguage*}{english}
\noindent This thesis investigated [problem] under [N] articulated
perspectives, adopting [approach]. The objectives were: (i) [obj1];
(ii) [obj2]; (iii) [obj3]. Research was conducted between [period],
with [N] articulated studies. Findings indicated [main result],
contributing to [field] through [contribution]. Limitations include
[limitation], with [agenda] suggested.

\noindent\textbf{Keywords}: term 1; term 2; term 3; term 4; term 5;
term 6.
\end{otherlanguage*}\end{resumo}

\pdfbookmark[0]{\listfigurename}{lof}
\listoffigures*\cleardoublepage
\pdfbookmark[0]{\listtablename}{lot}
\listoftables*\cleardoublepage

\begin{siglas}
\item[ABNT] Associação Brasileira de Normas Técnicas
\item[CAPES] Coordenação de Aperfeiçoamento de Pessoal de Nível Superior
\item[CEP] Comitê de Ética em Pesquisa
\item[CNPq] Conselho Nacional de Desenvolvimento Científico e Tecnológico
\item[FAP] Fundação de Amparo à Pesquisa
\item[ORCID] Open Researcher and Contributor ID
\item[PPG] Programa de Pós-Graduação
\item[TCLE] Termo de Consentimento Livre e Esclarecido
\end{siglas}

\pdfbookmark[0]{\contentsname}{toc}\tableofcontents*\cleardoublepage

\textual

% ═══════════════════════════════════════════════════
% MODALIDADE TRADICIONAL — comentar a partir daqui
% e descomentar a modalidade "por artigos" abaixo
% ═══════════════════════════════════════════════════

\chapter{INTRODUÇÃO GERAL}
\chapter{REFERENCIAL TEÓRICO}
\chapter{METODOLOGIA}
\chapter{ESTUDO 1: [Título]}
\chapter{ESTUDO 2: [Título]}
\chapter{ESTUDO 3: [Título]}
\chapter{DISCUSSÃO GERAL}
\chapter{CONSIDERAÇÕES FINAIS}

% ═══════════════════════════════════════════════════
% MODALIDADE TESE POR ARTIGOS — descomentar se for usar
% ═══════════════════════════════════════════════════
%
% \chapter{INTRODUÇÃO GERAL}
% \chapter{REFERENCIAL TEÓRICO COMUM}
% \chapter{METODOLOGIA GERAL}
%
% \chapter{ARTIGO 1: [Título]}
% \section{Status}
% [Publicado em [Revista, Ano] / Aceito / Submetido]
% % Para incluir PDF do artigo publicado:
% % \includepdf[pages=-]{artigos/artigo1.pdf}
%
% \chapter{ARTIGO 2: [Título]}
% \chapter{ARTIGO 3: [Título]}
%
% \chapter{DISCUSSÃO GERAL E CONSIDERAÇÕES FINAIS}

\postextual

\bibliography{referencias}

\begin{apendicesenv}
\partapendices
\chapter{INSTRUMENTOS DE COLETA}
\chapter{TCLE}
\end{apendicesenv}

\begin{anexosenv}
\partanexos
\chapter{PARECER DO CEP}
\chapter{COMPROVANTES DE PUBLICAÇÃO DOS ARTIGOS}
\end{anexosenv}

\end{document}
